\section*{Tests}
\label{sec:tests}

The main evaluation artifacts for this milestone are the automated test modules, the shared test infrastructure, and the \texttt{.tl} fixture programs.

\subsection*{Automated test modules}

\begin{longtable}{p{0.28\textwidth} p{0.66\textwidth}}
\toprule
\textbf{File} & \textbf{Purpose} \\
\midrule
\texttt{tests/test\_default.py} & Checks properties in \texttt{default} mode, where the entire state starts from a concrete zero initialization. These tests focus on whether the verifier can prove or refute invariants from a fixed start state. \\
\texttt{tests/test\_universal.py} & Checks properties in \texttt{universal} mode, where user variables and key turtle components are unconstrained initially. These tests validate reasoning under arbitrary initial conditions. \\
\texttt{tests/test\_specific.py} & Checks properties in \texttt{specific} mode using explicit parameter dictionaries. These tests validate the part of the pipeline that reads concrete input assignments from the command line and injects them into the start state. \\
\texttt{tests/test\_safety\_range.py} & Checks properties in \texttt{safety-range} mode. These tests validate the ghost-parameter mechanism and check whether the implementation can recover a symbolic safe region over initial inputs. \\
\texttt{tests/helpers.py} & Shared testing infrastructure. It builds the fixed-point system for a given program, constructs the property-evaluation context, suppresses verbose pipeline output during testing, and provides helper assertions for expected pass/fail outcomes. \\
\texttt{tests/run\_all.py} & Top-level test runner that discovers all \texttt{test\_*.py} files and executes the complete suite. \\
\bottomrule
\end{longtable}

\subsection*{Fixture programs used by the tests}

\begin{longtable}{p{0.24\textwidth} p{0.70\textwidth}}
\toprule
\textbf{Program file} & \textbf{Purpose in evaluation} \\
\midrule
\texttt{assign\_basic.tl} & Basic sequential assignments; used to test simple arithmetic invariants and no-motion/no-pen behavior. \\
\texttt{assign\_algebra.tl} & Multi-step arithmetic including multiplication and division; used to test algebraic propagation and nonlinear cases. \\
\texttt{conditional.tl} & Simple conditional branch; used to test branching semantics. \\
\texttt{loop\_basic.tl} & Counted loop with accumulation; used to test loop invariants. \\
\texttt{loop\_accum.tl} & Repetition-driven accumulator growth; used to test monotonic arithmetic invariants. \\
\texttt{loop\_cond.tl} & Loop with internal branching and arithmetic updates; used to test interaction between repetition and conditionals. \\
\texttt{goto\_computed.tl} & Computed \texttt{goto} destinations from symbolic variables; used for geometric safety properties. \\
\texttt{square\_goto.tl} & Multiple \texttt{goto} commands that trace a square; used to test bounded-region properties. \\
\texttt{forward\_square.tl} & Motion using \texttt{forward} and \texttt{right}; used for trig-dependent geometric tests. \\
\texttt{turns\_only.tl} & Pure heading updates; used for directional reasoning and normalization behavior. \\
\texttt{pen\_only.tl} & Pen-state updates without arithmetic or motion; used to validate \texttt{pendown}/\texttt{penup} semantics. \\
\texttt{pen\_with\_var.tl} & Pen commands combined with variable updates; used to test mixed symbolic and boolean state. \\
\texttt{pen\_toggle.tl} & Alternating pen state changes; useful for repeated pen-state transitions. \\
\texttt{param\_goto.tl} & Parameterized \texttt{goto}; used in \texttt{specific} and \texttt{safety-range} modes. \\
\texttt{param\_loop.tl} & Loop whose initial value depends on a parameter; used to test parameter-sensitive arithmetic invariants. \\
\texttt{param\_cond.tl} & Branching behavior controlled by a parameter; used to test mode-dependent initialization and branch selection. \\
\texttt{param\_scale.tl} & Arithmetic scaling with an input parameter; used to test concrete parameter injection and sign-sensitive properties. \\
\texttt{param\_pen.tl} & Parameterized arithmetic followed by pen commands; used to test the interaction of user inputs and pen state. \\
\texttt{read\_before\_write.tl} & Reads an input parameter before any local assignment; used to validate handling of externally supplied variables in \texttt{specific} mode. \\
\bottomrule
\end{longtable}

\subsection*{Individual test-case catalogue}

The suite contains 89 unit tests in total, grouped into four mode-specific modules.
Seven tests are automatically \emph{skipped} at runtime; all involving
programs that execute \texttt{right} or \texttt{forward} commands, producing
heading-dependent constraints that are too complex for \textsc{Spacer} to resolve
efficiently:

\begin{enumerate}[leftmargin=1.5em]
    \item \textbf{Heading normalization.}
    Each \texttt{right}~$\theta$ command updates the heading via a chain of
    nested \texttt{If}-expressions that wraps the result into $[0,360)$.
    The current encoding produces roughly 20 nested Z3 \texttt{If}-nodes per
    turn, causing the \textsc{Spacer} engine's invariant-inference loop to
    diverge or time out rather than converge on a stable inductive
    strengthening.

    \item \textbf{Trigonometric motion.}
    \texttt{forward}~$d$ updates position as
    $\mathit{xcor}' = \mathit{xcor} + d\cdot\sin(\mathit{heading})$ and
    $\mathit{ycor}' = \mathit{ycor} + d\cdot\cos(\mathit{heading})$.
    Because SMT solvers do not natively support transcendental functions, the
    encoding approximates $\sin$ and $\cos$ with interval
    over-approximations over discretised heading buckets.  These
    approximations further inflate the constraint size and interact
    multiplicatively with the heading-normalization chain.
\end{enumerate}

\noindent
Crucially, the skips reflect a \emph{solver scalability} limitation, not a
defect in the encoding logic.  The affected programs are correctly translated
into CHC rules, and the properties are well-formed; \textsc{Spacer} simply
cannot synthesize a sufficiently strong inductive invariant within the
configured timeout when both heading normalization and trigonometric
approximations are present in the same fixed-point system.  Programs that
modify heading but do not use \texttt{forward}/\texttt{backward} (and thus
avoid the trig layer) are handled correctly, as demonstrated by the four
passed directional tests in default mode and the four active directional
tests in safety-range mode.  Similarly, programs using \texttt{goto} (which
sets position directly without trigonometry) pass all geometric tests across
all four modes.

\subsection*{Running the tests}

This suite requires only the \texttt{unittest} module from the Python standard library, other than the dependecies for CHC Verfication and Chiron, which is available in any standard Python 3 installation.

The full automated suite can be run with:

\begin{lstlisting}[style=terminal, language=bash]
cd Chiron-Framework/ChironCore/CHC_Verification/tests
python run_all.py
\end{lstlisting}

\noindent Narrower invocations are also supported:

\begin{lstlisting}[style=terminal, language=bash]
python run_all.py -v                        # verbose
python -m unittest test_default             # one module
python -m unittest test_default.TestDefaultArithmetic.test_arith_assign_z_nonneg_pass
\end{lstlisting}
